\documentclass[10pt,a4paper]{article}
\usepackage[T1]{fontenc}
\usepackage[utf8]{inputenc}
\usepackage{lmodern}
\usepackage{amsmath,amssymb}
\usepackage[margin=21mm]{geometry}
\usepackage{microtype}
\usepackage[hidelinks]{hyperref}
\usepackage{enumitem}
\setlist[itemize]{leftmargin=*,nosep,topsep=3pt}
\setlength{\parindent}{0pt}
\setlength{\parskip}{5pt}
\setlength{\emergencystretch}{3em}
\newcommand{\question}[1]{\par\medskip{\normalsize\bfseries #1}\par\smallskip}
\pagestyle{plain}

\begin{document}
{\Large\bfseries Synaptic Changes in Pallidostriatal Circuits Observed in the Parkinsonian Model Triggers Abnormal Beta Synchrony with Accurate Spatio-temporal Properties across the Basal Ganglia\par}

\textbf{Paper:} Azizpour Lindi, Mallet and Leblois (2024).\quad Experimentally constrained rate and spiking models.\par
\href{https://doi.org/10.1523/JNEUROSCI.0419-23.2023}{DOI: \nolinkurl{10.1523/JNEUROSCI.0419-23.2023}}\par
{\footnotesize Page references below use PDF page numbers, including front matter.}

\question{To which general question does this paper relate?}
Mechanisms of beta synchronization and model structure. Tests whether experimentally constrained loops reproduce both frequency and relative spike timing.

\question{How does it answer the general questions?}
\textbf{Finding.} The model includes STN, prototypic/arkypallidal GPe, striatal fast-spiking interneurons and D2 medium spiny neurons. Strong STN--GPe or intra-GPe feedback favours gamma; pallidostriatal loops favour lower beta. Combined synaptic changes reproduce parkinsonian beta and phase relationships (PDF pp.~1--5, 8--16).

\textbf{Usable insight.} In model perturbations, GPe inhibition abolishes oscillations whereas STN inhibition does not, matching cited rat experiments. Reproducing phase relationships is an additional model-validation criterion beyond matching a spectral peak. These results explain synchronization, not a demonstrated clinical benefit of suppressing it.

\textbf{Model relevance.} The rate approximation uses delayed, synaptically filtered interactions and threshold-linear activation; synaptic time constants differ by projection. The notebook can represent an augmented rate approximation with extra filter states, but the spiking model is not exactly its memoryless rate equation. Connectivity, membrane, synaptic and delay parameters are tabulated in Methods (PDF pp.~3--6). A new reduction and validated observation/actuator maps are needed.

\question{What papers does it cite that could provide more information?}
\begin{itemize}
\item De la Crompe et al. (2020), \emph{The globus pallidus orchestrates abnormal network dynamics in a model of parkinsonism}: perturbation constraints (PDF p.~18).
\item Corbit et al. (2016), \emph{Pallidostriatal projections promote beta oscillations in a dopamine-depleted biophysical network model}: related loop hypothesis (PDF p.~18).
\end{itemize}

\end{document}
